\chapter{Domain and Technical Foundations}
\label{ch:foundations}

To contextualize the Wolverine prototype, this chapter outlines the foundational concepts in cross-platform identity intelligence, entity resolution, and the specific data structures utilized within the system's simulated ecosystem.

\section{Relevant Domain Concepts}
\label{sec:foundations-domain}
Modern intelligence gathering increasingly relies on analyzing digital artifacts scattered across independent domains. In the context of the Wolverine prototype, the primary domain is modeled after Open Source Intelligence (OSINT) and dark web monitoring. 
Key concepts include:
\begin{itemize}
    \item \textbf{Heterogeneous Silos:} Independent digital communities (forums, marketplaces, social networks) that enforce distinct schemas and share no underlying infrastructure.
    \item \textbf{Digital Footprints:} The trail of metadata left by a user, including aliases, timestamps, IP addresses (or proxies), and interaction patterns.
    \item \textbf{Cross-Platform Correlation:} The analytical process of deducing that two distinct digital personas on different platforms belong to the same real-world entity.
\end{itemize}

\section{Data and Entity Representation}
\label{sec:foundations-data-rep}
A significant challenge in correlating heterogeneous data is structural incompatibility. Platform A might represent a user interaction as a nested JSON object involving \texttt{buyer\_id} and \texttt{seller\_id}, whereas Platform B uses a flat relational schema of \texttt{author} and \texttt{post}.

Wolverine addresses this by enforcing a \textit{Canonical Record Schema} during the ingestion pipeline. Raw captures are mapped into a standardized format containing:
\begin{itemize}
    \item \textbf{Entity Nodes:} Representations of personas, digital artifacts (e.g., cryptocurrency wallets), or content.
    \item \textbf{Relational Edges:} Defined interactions (e.g., \texttt{AUTHORED}, \texttt{INTERACTED\_WITH}, \texttt{TRANSACTED}).
    \item \textbf{Temporal Metadata:} Standardized ISO-8601 timestamps used to evaluate proximity constraints.
\end{itemize}
This canonical representation is the foundational layer upon which the heuristic resolution algorithms operate.

\section{Graph Concepts in Wolverine}
\label{sec:foundations-graph}
The canonical data model naturally forms a directed graph where personas and artifacts are nodes, and their activities form the connecting edges. 
In theoretical production deployments, resolving identities across such a network requires highly scalable graph traversal methodologies, often leveraging specialized graph databases or recursive Common Table Expressions (CTEs) in relational systems like PostgreSQL.

However, as a prototype, Wolverine implements a computationally bounded approach:
\begin{itemize}
    \item \textbf{In-Memory Breadth-First Search (BFS):} Graph projection and neighborhood traversal are executed entirely in the memory space of the Node.js application process. 
    \item \textbf{Bounded Depth:} Traversals are strictly limited to a shallow depth (typically 2 to 4 hops) to prevent memory exhaustion and infinite loops in cyclic interaction patterns.
\end{itemize}
This foundation enables the demonstration of multi-hop intelligence correlation while deliberately constraining the operational footprint to a single host.

\section{Relevant Computational Methods}
\label{sec:foundations-methods}
To determine if two distinct aliases might represent the same entity, Wolverine employs deterministic heuristic similarity scoring. Rather than relying on computationally expensive or unpredictable Machine Learning (ML) models, the prototype uses the following formal methods:

\subsection{Jaro-Winkler Distance}
The Jaro-Winkler string metric is a measure of edit distance between two sequences. It is particularly effective for short strings such as usernames or aliases, as it heavily weights prefix matches. The Jaro-Winkler similarity $s_{jw}$ is bounded between 0 and 1, where 1 equates to an exact match. Wolverine uses this metric as the primary driver for computing `aliasSimilarity`.

\subsection{Temporal Proximity Decay}
Behavioral correlation assumes that coordinated actions across platforms often occur within narrow time windows. Wolverine applies a linear decay function to calculate a temporal proximity score based on the day difference between account creations or specific activities. As the temporal distance increases, the proximity score linearly decays toward zero.

\subsection{Heuristic Weighting}
The final resolution score is a weighted linear combination of distinct similarity features (alias, display name, temporal proximity, and behavioral overlap). By applying static, predefined weights to these normalized features, the system computes a deterministic confidence score that governs whether two nodes should be linked in the graph. 
\textit{Detailed formulations of these models, including prototype-specific static constants, are provided in \Cref{ch:analytical-methodology}.}
